\documentclass[tikz,border=10pt]{standalone}
\usepackage[no-math]{fontspec}
\usepackage[brazil]{babel}
\usepackage{amsmath,amssymb}
\IfFileExists{C:/Windows/Fonts/Lato_Regular.ttf}{%
  \setmainfont{Lato}[
    Path = C:/Windows/Fonts/,
    UprightFont = Lato_Regular.ttf,
    BoldFont = Lato_Bold.ttf,
    ItalicFont = Lato_Italic.ttf,
    BoldItalicFont = Lato_BoldItalic.ttf
  ]
  \setsansfont{Lato}[
    Path = C:/Windows/Fonts/,
    UprightFont = Lato_Regular.ttf,
    BoldFont = Lato_Bold.ttf,
    ItalicFont = Lato_Italic.ttf,
    BoldItalicFont = Lato_BoldItalic.ttf
  ]
}{%
  \RequirePackage[default]{lato}
}
\renewcommand{\familydefault}{\sfdefault}

\usepackage{xcolor}
\usepackage{tikz}
\usepackage{pgfplots}
\pgfplotsset{compat=1.18}
\usetikzlibrary{arrows.meta,positioning,shapes.geometric,fit,backgrounds,calc,decorations.pathreplacing,patterns}

% Paleta institucional UEFS / Sociologia Rural
\definecolor{srNavy}{HTML}{1F4268}
\definecolor{srNavyDark}{HTML}{152E4A}
\definecolor{srGold}{HTML}{B38A3D}
\definecolor{srGoldDark}{HTML}{8C6927}
\definecolor{srTeal}{HTML}{0E7C7B}
\definecolor{srForest}{HTML}{2D6A4F}
\definecolor{srAlert}{HTML}{9E2A2B}
\definecolor{srInk}{HTML}{26313D}
\definecolor{srMuted}{HTML}{5B6573}

\definecolor{bgNavy}{HTML}{EDF2F7}
\definecolor{bgGold}{HTML}{F8F3E9}
\definecolor{bgTeal}{HTML}{EBF6F5}
\definecolor{bgForest}{HTML}{EEF5F0}
\definecolor{bgAlert}{HTML}{FDF2F2}

\tikzset{
  >=Stealth,
  srBox/.style={
    draw=#1, fill=white, rounded corners=5pt, line width=1.15pt,
    inner sep=7pt, align=center, font=\footnotesize\sffamily, text=srInk
  },
  srTint/.style={
    draw=#1, fill=#1!8!white, rounded corners=5pt, line width=1.15pt,
    inner sep=7pt, align=center, font=\footnotesize\sffamily, text=srInk
  },
  srBadge/.style={
    fill=#1, text=white, rounded corners=9pt, inner xsep=8pt, inner ysep=3pt,
    font=\scriptsize\bfseries\sffamily
  },
  srTitle/.style={font=\normalsize\bfseries\sffamily, text=srNavy, align=center},
  srSub/.style={font=\footnotesize\sffamily, text=srMuted, align=center},
  srArrow/.style={->, >=Stealth, line width=1.35pt},
  srDArrow/.style={<->, >=Stealth, line width=1.35pt}
}


\begin{document}
\begin{tikzpicture}

  % Title
  \node[srTitle] (title) at (0, 4.45) {Organização Social da Agricultura Familiar (Chayanov e Lei nº 11.326/2006)};
  \node[srSub] (sub) at (0, 3.98) {Contraste entre racionalidades produtivas e critérios legais de enquadramento no Brasil};

  % Top Left Card: Lógica Camponesa
  \node[srTint=srTeal, text width=6.1cm, minimum height=4.25cm, anchor=north] (left) at (-4.25, 3.30) {};
  \node[srBadge=srTeal, rounded corners=5pt, anchor=north] at ([yshift=-9pt]left.north) {Lógica Camponesa (Chayanov)};
  \node[font=\footnotesize\sffamily, text=srInk, align=left, text width=5.6cm, anchor=north] at ([yshift=-36pt]left.north) {%
    \textbullet\ \textbf{Unidade produção = consumo}: família e estabelecimento indivisíveis\\[6pt]
    \textbullet\ \textbf{Balanço trabalho--penosidade}: equilíbrio entre esforço e necessidade\\[6pt]
    \textbullet\ \textbf{Pluriatividade e autoconsumo}: reprodução social e segurança alimentar
  };

  % Top Right Card: Empresa Patronal Capitalista
  \node[srTint=srGold, text width=6.1cm, minimum height=4.25cm, anchor=north] (right) at (4.25, 3.30) {};
  \node[srBadge=srGoldDark, rounded corners=5pt, anchor=north] at ([yshift=-9pt]right.north) {Empresa Patronal Capitalista};
  \node[font=\footnotesize\sffamily, text=srInk, align=left, text width=5.6cm, anchor=north] at ([yshift=-36pt]right.north) {%
    \textbullet\ \textbf{Separação capital/trabalho}: gestão dissociada da execução direta\\[6pt]
    \textbullet\ \textbf{Lucro e taxa de retorno}: acumulação orientada pelo mercado\\[6pt]
    \textbullet\ \textbf{Trabalho assalariado predominante}: contratação externa contínua ou sazonal
  };

  % Comparison connector
  \draw[srDArrow, srMuted] (left.east) -- (right.west)
    node[midway, fill=white, inner sep=3pt, font=\scriptsize\bfseries\sffamily, text=srMuted] {VS.};

  % Bottom Unified Card: 4 Requisitos Legais
  \node[srTint=srForest, text width=14.6cm, minimum height=3.75cm, anchor=north] (bottom) at (0, -1.85) {};
  \node[srBadge=srForest, rounded corners=5pt, anchor=north] at ([yshift=-9pt]bottom.north) {4 Requisitos Legais da Agricultura Familiar (Lei nº 11.326/2006)};

  % 4 Inner pills (2x2 grid)
  \node[srBox=srForest, text width=6.4cm, minimum height=0.88cm, align=left, anchor=north east] (p1)
    at ([xshift=-0.25cm, yshift=-38pt]bottom.north)
    {\textbf{(1)} Área $\le$ 4 módulos fiscais};

  \node[srBox=srForest, text width=6.4cm, minimum height=0.88cm, align=left, anchor=north west] (p2)
    at ([xshift=0.25cm, yshift=-38pt]bottom.north)
    {\textbf{(2)} Mão de obra predominantemente familiar};

  \node[srBox=srForest, text width=6.4cm, minimum height=0.88cm, align=left, anchor=north] (p3)
    at ([yshift=-0.24cm]p1.south)
    {\textbf{(3)} Renda oriunda do próprio estabelecimento};

  \node[srBox=srForest, text width=6.4cm, minimum height=0.88cm, align=left, anchor=north] (p4)
    at ([yshift=-0.24cm]p2.south)
    {\textbf{(4)} Gestão dirigida pela família};

  % Downward arrow from Lógica Camponesa to Lei 11.326
  \draw[srArrow, srTeal] (left.south) -- ++(0,-0.42) -| (bottom.north);

\end{tikzpicture}
\end{document}
